\section*{Advanced Decision Cases: Organization Design When Incentives Collide}
\addcontentsline{toc}{section}{Advanced Decision Cases: Organization Design When Incentives Collide}

\subsection*{Case A. A High-Performing Salesperson Sells an Impossible Course}
\addcontentsline{toc}{subsection}{Case A. A High-Performing Salesperson Sells an Impossible Course}

Assume a small educational firm rewards admissions staff primarily for paid enrollments. A salesperson discovers that promising individualized 1-to-1 guidance improves conversion and begins describing the program that way. The standard offer, however, includes limited group mentorship and 2 scheduled project reviews. The salesperson brings in 15 additional learners at ₹25,000 each, apparently adding ₹3,75,000 of revenue. The founder initially celebrates the achievement.

The real obligation is different. Suppose each learner reasonably expects 8 additional hours of skilled individual attention. The new promise implies $15 \times 8 = 120$ hours of work. At a normalized rate of ₹1,500 per qualified expert hour, that is ₹1,80,000 before scheduling, support, customer acquisition, and rework. More seriously, the firm's qualified capacity may simply not exist; the financial estimate understates the risk. Teachers must either break the promise, reduce service to other learners, or exhaust themselves. The sales team has optimized its metric by exporting cost and reputational damage to delivery.

This is a design failure across 4 mechanisms:
\begin{enumerate}
    \item \textbf{Incentives:} Reward is tied to gross enrollment without promise quality.
    \item \textbf{Decision rights:} The salesperson may alter service claims without approval.
    \item \textbf{Information:} Sales lacks live visibility into review capacity.
    \item \textbf{Culture:} Management rewards the numbers before validating their legitimacy.
\end{enumerate}

Simply firing the employee or demanding ``more integrity'' ignores why the behavior was attractive and tolerated. A repaired design fixes the authorized offer, publishes the exact service boundary, requires exceptions to be approved before payment, connects the sales reward to valid contracts and appropriate realization, and gives the operations owner a credible veto over infeasible commitments. The founder must also rectify promises already made and protect the affected learners. The deeper lesson is that every sale designs part of the future organization: when commercial and operating rights are disconnected, more revenue can mean less institutional value.

\subsection*{Case B. 2 Excellent Managers Produce 1 Failing Customer Project}
\addcontentsline{toc}{subsection}{Case B. 2 Excellent Managers Produce 1 Failing Customer Project}

Imagine a college partnership led jointly by an academic program manager and a commercial account manager. Both are competent. The academic manager promises rigorous assessment dates but cannot get confirmed classroom access; the account manager promises the college a revised deadline without consulting teaching; finance is not informed that the final invoice now depends on a changed acceptance document. Each manager can explain their individual decision. The project as a whole fails.

A conventional reaction is to appoint another senior manager. That may merely add a 3rd person to the confusion. The real missing element is an end-to-end project owner and clear decision rights over scope, calendar, acceptance, and escalation. A decision table can assign the academic lead authority over instructional standards, the account manager over ordinary customer communication, the project owner over integrated schedule and accepted scope, and the founder over material financial or ethical exceptions. The process should specify which decisions require mandatory consultation and who records the agreed plan.

The new structure must also solve information availability. A shared milestone record shows classroom confirmation, delivery readiness, reports due, invoice triggers, and decision owners. Weekly exception review focuses on unmet dependencies. When a college requests a major shift, the account manager can explain which decisions need review rather than offering an informal commitment that binds the organization without authority. This illustrates why \textbf{competence is not a substitute for coordination design}. A team can hire excellent people and still be a poor organization. The performance test is whether multiple correct local decisions combine into a correct customer result.

\subsection*{Case C. A Faculty Star Weakens the Faculty Institution}
\addcontentsline{toc}{subsection}{Case C. A Faculty Star Weakens the Faculty Institution}

A future technical institute recruits a highly regarded educator whose lectures attract demand. The faculty member also insists on exclusive control of grading, declines calibration, and refuses to share instructional resources. Their learners give excellent presentation ratings, but independent assessments reveal inconsistent standards across courses. Leadership hesitates to intervene because the educator is commercially valuable.

The conflict concerns professional autonomy versus institutional accountability. Academic freedom and teaching discretion are not equivalent to a right to change evaluation standards without scrutiny. The organization should preserve the educator's legitimate methods while defining shared requirements for outcomes, integrity, fairness, and evidence. Peer calibration can examine actual student work rather than demanding conformity of teaching style. The institute can also recognize the faculty member's intellectual contribution through attribution and agreed rights over materials.

A decision-rights framework separates what the educator can decide independently (examples, instructional activities, pacing within approved boundaries) from what requires shared academic governance (pass thresholds, high-stakes credential claims, integrity sanctions, and modifications affecting cohort comparability). Compensation and promotion should reflect institutional contribution, not only demand generated by a celebrated lecturer.

The short-run risk is losing a prominent person. The long-run risk is becoming a collection of personal brands with incompatible standards. A durable academic institution must be able to respect expertise while governing the claims it makes under its own name. If standards are only enforceable against junior people, the culture teaches that prestige outranks integrity.

\subsection*{Case D. Founder Delegation Fails Because the Founder Keeps Deciding}
\addcontentsline{toc}{subsection}{Case D. Founder Delegation Fails Because the Founder Keeps Deciding}

Consider a founder who hires a capable operations coordinator and announces that the coordinator owns learner onboarding. In practice, the founder personally answers most learner messages, changes deadlines in private conversations, and overrides the coordinator's scheduling without updating the shared system. The coordinator becomes reluctant to decide because their decisions may be contradicted. The founder concludes that the new hire lacks initiative and resumes control.

The failure is not proof that delegation was premature; it may reflect contradictory authority. The repair begins with a delegation contract. The coordinator may make ordinary schedule adjustments within explicit policy, approve specified low-cost service remedies, request missing intake documents, and escalate only contractual, academic, financial, or safety exceptions. All customer-facing decisions are recorded in 1 authoritative place. The founder agrees not to make parallel private commitments and, when an override is necessary, to explain it through the same workflow.

Then run a bounded test: for 4 weeks, the coordinator owns 20 representative onboarding cases. Measure outcome quality, missed commitments, escalation frequency, learner confusion, and founder intervention time. A small number of justified escalations is not a failure. Success means the coordinator handles the normal case competently and recognizes the abnormal case early. The founder should also inspect whether instructions were clear and resources sufficient before evaluating the person. The organizational lesson is uncomfortable: a founder can be the reason their organization remains founder-dependent. Delegation requires surrendering the psychological reward of being indispensable and allowing other competent people to exercise bounded discretion.

\subsection*{Case E. Change Is Announced but the Incentives Remain Unchanged}
\addcontentsline{toc}{subsection}{Case E. Change Is Announced but the Incentives Remain Unchanged}

Suppose an AI consulting firm announces a new commitment to reusable products and reduced customization. Leaders ask consultants to document reusable code, reject low-value scope changes, and invest time in shared components. Yet performance reviews still reward only billed utilization and current-quarter revenue. Staff naturally devote attention to billable client work, while reusable-asset creation becomes unpaid evening labor.

The strategy has changed rhetorically, but the organizational design has not. Galbraith's logic suggests examining all 5 levers: strategy, structure, processes, rewards, and people. The firm may need protected time for asset development, a product owner to prioritize the shared library, transparent rules for accepting customization, incentives recognizing verified reuse, and training in modular engineering. Otherwise the new aspiration competes with the old reward structure and almost certainly loses in daily decisions.

A test might allocate 10\% of a small team's capacity for 6 weeks to building a reusable evaluation harness for common client workflows. Measure how often it is adopted by different projects, actual setup time saved, defect reduction, maintenance burden, and lost billable opportunity. If the asset is seldom reused, the team should not praise the documentation volume; it should revisit whether the portfolio contains a sufficiently common problem. The design must respond to evidence. The lesson is that \textbf{organizational change requires an altered pattern of decisions and resource commitments}, not an altered vocabulary. Culture follows what people see rewarded, enabled, and protected under pressure.

\subsection*{Case F. Succession Without Transfer of Legitimacy}
\addcontentsline{toc}{subsection}{Case F. Succession Without Transfer of Legitimacy}

Imagine a mature educational institute whose founder appoints a new chief executive but continues to negotiate major contracts, approve faculty promotions, and settle disputes through informal messages. Donors and staff still approach the founder as the real authority. The successor has the title but cannot credibly commit the institution. Board meetings appear orderly while the informal organization contradicts the formal 1.

Succession must transfer 3 different things: \textbf{decision rights}, \textbf{institutional knowledge}, and \textbf{stakeholder legitimacy}. The board or lawful governance body should define the successor's authority and the founder's remaining role, with clear procedures for exceptional consultation. Key external stakeholders should hear consistent messages. Critical contracts, academic decisions, and personnel matters should route through the new authority. If the founder retains a governance role, its boundaries should be explicit and respected.

A transition scorecard can track decisions made by the new leader without reversal, clarity of external counterparties, retention and development of key staff, financial and academic performance, and whether the founder has stopped serving as the routine escalation point. Good succession is not a test of whether the successor imitates the founder. It is a test of whether the institution can preserve its purpose through legitimately different leadership.

\subsection*{A Paired Architecture Test: When Do You Need an Organization Instead of Another Tool?}
\addcontentsline{toc}{subsection}{A Paired Architecture Test: When Do You Need an Organization Instead of Another Tool?}

A founder observing delivery delays often has 3 apparent remedies: buy software, recruit a specialist, or reorganize the team. These options operate at different causal levels. Software can reduce the effort of a task with known inputs and outputs. A specialist can add judgment or capacity. Organization design can clarify authority, incentives, and the flow of decisions. None automatically substitutes for the other.

Consider a hypothetical backlog of college-partner reporting. Reports take 4 hours of analyst labor and another 10 days for approvals. Buying an AI report generator may reduce drafting from 4 hours to 1, yet cash collection remains delayed because nobody knows who may approve the final report. Hiring an additional analyst increases drafting capacity but can make the approval queue worse. A simple decision-rights agreement with the college (authorized approver, acceptance standard, response interval, and escalation) may reduce most waiting time. Only after that change does better drafting technology create a large additional benefit.

The test is to identify where time, errors, and uncertainty actually accumulate; separate production bottlenecks from governance bottlenecks; then choose the intervention with the best expected reduction in total customer-impacting cost. The founder should not invest in a sophisticated answer to a problem they have not diagnosed. Conversely, if the approval process is clear but too many legitimate reports require hours of expert analysis, capacity or automation might indeed be the right answer.

A disciplined team records a counterfactual: what would happen if we changed only the tool, only the role, or only the authority? It need not claim precise causal certainty. It needs a credible theory that connects the chosen intervention to a measurable customer outcome. That is where operational excellence and organization design converge: \textbf{the quality of the enterprise depends on the quality of its diagnosis before it commits resources}.

\subsection{A Closing Synthesis: An Institution Outlives Its Central Personality}
\label{sec:closing_synthesis_institution}

The earliest company may be indistinguishable from the founder's attention. A mature institution must become more than that attention: qualified people can act, disagree, teach, improve, govern, and replace 1 another while preserving the purpose that brought the organization into existence. This requires aligned structure, process, incentives, talent, culture, and decision rights, not merely scale in headcount.

Bayes Institute's long-term ambition makes the standard unusually demanding. An educational institution devoted to human possibility should itself be capable of developing human potential inside its organization. Teachers must become better teachers, researchers more rigorous researchers, operators more capable owners, and leaders better at entrusting others. \textbf{The founder succeeds institutionally when the mission becomes less dependent on the founder's individual judgment and more reliably embodied in the collective system.}

\begin{quote}
\textit{The purpose of organization design is not to give the founder more people to command. It is to build an arrangement in which competent people can make consequential decisions together, uphold the mission under pressure, and continuously create capability beyond the founder's personal reach.}
\end{quote}

\section*{Research and Intellectual Foundations}
\addcontentsline{toc}{section}{Research and Intellectual Foundations}

These references identify the foundational research and scholarship on organizational design, culture, incentives, and institutional evolution.

\begin{enumerate}
    \item[\textbf{[G1]}] Galbraith, J. R., \textit{The Star Model} (strategy, structure, processes, rewards, people), Galbraith Management Consultants.
    \item[\textbf{[G2]}] Mintzberg, H. (1980), ``Structure in 5's: A Synthesis of the Research on Organization Design,'' \textit{Management Science}, 26(3), pp. 322-341.
    \item[\textbf{[G3]}] Bain \& Company, \textit{The 5 Steps to Better Decisions} and the RAPID decision-rights framework.
    \item[\textbf{[G4]}] Holmstr\"{o}m, B., \& Milgrom, P. (1991), ``Multitask Principal-Agent Analyses: Incentive Contracts, Asset Ownership, and Job Design,'' \textit{Journal of Law, Economics, and Organization}, 7, pp. 24-52.
    \item[\textbf{[G5]}] Schein, E. H., \& Schein, P. (2016), \textit{Organizational Culture and Leadership}, 5th edition, Wiley.
    \item[\textbf{[G6]}] Edmondson, A. (1999), ``Psychological Safety and Learning Behavior in Work Teams,'' \textit{Administrative Science Quarterly}, 44(2), pp. 350-383.
    \item[\textbf{[G7]}] Kotter, J. P. (1996), \textit{Leading Change}, Harvard Business School Press.
    \item[\textbf{[G8]}] Hughes, M. (2011), ``Do 70 Per Cent of All Organizational Change Initiatives Really Fail?'' \textit{Journal of Change Management}, 11(4), pp. 451-464.
    \item[\textbf{[G9]}] March, J. G. (1991), ``Exploration and Exploitation in Organizational Learning,'' \textit{Organization Science}, 2(1), pp. 71-87.
    \item[\textbf{[G10]}] Nonaka, I. (1994), ``A Dynamic Theory of Organizational Knowledge Creation,'' \textit{Organization Science}, 5(1), pp. 14-37.
    \item[\textbf{[G11]}] National Institute of Standards and Technology, \textit{Baldrige Excellence Framework} (leadership, strategy, customers, measurement, workforce, operations, results).
    \item[\textbf{[G12]}] Hackman, J. R. (2002), \textit{Leading Teams: Setting the Stage for Great Performances}, Harvard Business School Press.
\end{enumerate}

\subsubsection*{Further Reading}

Herbert A. Simon, \textit{Administrative Behavior}; Alfred D. Chandler Jr., \textit{Strategy and Structure}; Richard Rumelt, \textit{Good Strategy/Bad Strategy}; L. David Marquet, \textit{Turn the Ship Around!}; Michael C. Jensen and William H. Meckling on agency theory; Charles O'Reilly and Jeffrey Pfeffer on talent and culture. These authors offer useful contrasting lenses, not 1 unified doctrine.

\subsubsection*{Boundary to Part VII}

Chapters 23-26 have shown how to select people, evaluate specialized talent, design reliable processes, and arrange human authority. Chapter 27 turns to scaling: how these arrangements behave when the firm grows, what breaks under volume, and how management must preserve quality and unit economics while crossing organizational thresholds.
